\section{Introduction}


The design of a token economy for a multi-chain blockchain platform presents unique challenges that do not arise in single-chain systems. In a multi-chain architecture such as Lux Network, the native token must serve several distinct economic functions simultaneously: securing consensus through staking, pricing computational resources through transaction fees, coordinating economic activity across sovereign appchains, and providing a store of value for long-term network participants \cite{nakamoto2008bitcoin, buterin2014ethereum}.

Traditional tokenomic models, such as Bitcoin's fixed-supply halving schedule \cite{nakamoto2008bitcoin} or Ethereum's EIP-1559 base fee mechanism \cite{roughgarden2020eip1559}, are designed for single-chain environments and do not account for the cross-chain dynamics inherent in multi-appchain architectures. The Cosmos model \cite{kwon2016cosmos} addresses multi-chain economics through the ATOM staking token but relies on inflationary rewards that dilute non-staking holders. Polkadot's candle auction mechanism \cite{wood2016polkadot} ties parachain access to token lockup, creating capital inefficiency.

This paper introduces a unified tokenomic framework for Lux Network that addresses these limitations through three key innovations:

\begin{enumerate}[label=(\roman*)]
    \item \textbf{Adaptive reward functions}: Staking rewards adjust dynamically based on the ratio of staked to circulating supply, targeting an optimal staking ratio that balances network security with liquidity.
    \item \textbf{Multi-tier fee burning}: Transaction fees are partially burned at rates that vary by appchain type, creating deflationary pressure proportional to network utilization.
    \item \textbf{Appchain economic coupling}: A portion of appchain validator rewards derives from primary network activity, creating positive-sum economic incentives for ecosystem growth.
\end{enumerate}

The remainder of this paper is organized as follows. Section~\ref{sec:supply} defines the token supply model and distribution schedule. Section~\ref{sec:staking} formalizes the staking reward mechanism. Section~\ref{sec:fees} introduces the adaptive fee-burning protocol. Section~\ref{sec:game} presents game-theoretic analysis of validator behavior. Section~\ref{sec:appchain} discusses appchain economic coupling. Section~\ref{sec:simulation} provides simulation results. Section~\ref{sec:comparison} compares with existing systems. Section~\ref{sec:conclusion} concludes.

